\documentclass[12pt,a4paper,oneside]{report}
\usepackage[utf8]{inputenc}
\usepackage[T1]{fontenc}
\usepackage[french]{babel}
\usepackage{lmodern}
\usepackage{geometry}
\usepackage{setspace}
\usepackage{hyperref}
\usepackage{graphicx}
\usepackage{booktabs}
\usepackage{longtable}
\usepackage{array}
\usepackage{enumitem}
\usepackage{xcolor}
\usepackage{fancyhdr}
\usepackage{titlesec}
\usepackage{tocloft}
\usepackage{listings}
\usepackage{float}
\usepackage{caption}
\usepackage{pdflscape}
\usepackage{titletoc}

\geometry{margin=2.5cm}
\setstretch{1.23}
\setlength{\parindent}{0pt}
\setlength{\parskip}{0.55em}
\setlist[itemize]{itemsep=0.35em, topsep=0.45em}
\setlist[enumerate]{itemsep=0.35em, topsep=0.45em}

\definecolor{PrimaryBlue}{HTML}{1F4E79}
\definecolor{AccentBlue}{HTML}{2D6A9F}
\definecolor{SoftGray}{HTML}{F4F7FA}
\definecolor{BorderGray}{HTML}{D7DEE8}
\definecolor{TextGray}{HTML}{2E3440}
\hypersetup{
  colorlinks=true,
  linkcolor=PrimaryBlue,
  urlcolor=AccentBlue,
  pdftitle={Rapport technique - Plateforme de simulation de phishing et scoring de risque humain},
  pdfauthor={Equipe Projet}
}

\pagestyle{fancy}
\fancyhf{}
\fancyhead[L]{Rapport technique professionnel}
\fancyhead[R]{\leftmark}
\fancyfoot[C]{\thepage}
\renewcommand{\headrulewidth}{0.5pt}
\renewcommand{\footrulewidth}{0pt}

\titleformat{\chapter}[display]
  {\bfseries\Huge\color{PrimaryBlue}}
  {\filleft\fontsize{48}{52}\selectfont\color{PrimaryBlue!20}\thechapter}
  {-3.5ex}
  {\titlerule[1pt]\vspace{1ex}\filright}
  [\vspace{0.8ex}\titlerule]
\titleformat{\section}
  {\bfseries\Large\color{PrimaryBlue}}
  {\thesection}{0.7em}{}
\titleformat{\subsection}
  {\bfseries\large\color{AccentBlue}}
  {\thesubsection}{0.7em}{}
\titlespacing*{\chapter}{0pt}{-5pt}{18pt}
\titlespacing*{\section}{0pt}{14pt}{6pt}
\titlespacing*{\subsection}{0pt}{10pt}{4pt}

\setcounter{secnumdepth}{3}
\setcounter{tocdepth}{2}
\renewcommand{\cftchapfont}{\bfseries\color{PrimaryBlue}}
\renewcommand{\cftsecfont}{\color{TextGray}}
\renewcommand{\cftsubsecfont}{\color{TextGray}}
\renewcommand{\cftchappagefont}{\bfseries\color{PrimaryBlue}}
\renewcommand{\cftsecpagefont}{\color{TextGray}}

\lstset{
  basicstyle=\ttfamily\small,
  breaklines=true,
  frame=single,
  backgroundcolor=\color{SoftGray},
  rulecolor=\color{BorderGray},
  columns=fullflexible
}

\newcommand{\keypoint}[2]{
\vspace{0.6em}
\begin{center}
\fcolorbox{BorderGray}{SoftGray}{
\begin{minipage}{0.94\linewidth}
\textbf{\color{PrimaryBlue}#1}\par\smallskip
#2
\end{minipage}}
\end{center}
\vspace{0.4em}
}

\newcommand{\captureplaceholder}{
\begin{center}
\fcolorbox{BorderGray}{SoftGray}{
\begin{minipage}[c][6.2cm][c]{0.9\linewidth}
\centering
\textbf{\color{AccentBlue}[Inserer capture ici]}
\end{minipage}}
\end{center}
}

\begin{document}

\begin{titlepage}
  \centering
  {\Large\bfseries Rapport technique complet\par}
  \vspace{1cm}
  {\Huge\bfseries Plateforme de simulation de phishing\par}
  {\Huge\bfseries et scoring de risque humain\par}
  \vspace{1.2cm}
  {\Large Backend + Client (Frontend)\par}
  \vspace{1.5cm}
  {\large Document de type rapport de stage / fin d'etudes\par}
  \vfill
  {\large Annee universitaire 2025-2026\par}
  {\large Date: \today\par}
\end{titlepage}

\tableofcontents
\listoftables
\clearpage

\chapter*{Resume}
\addcontentsline{toc}{chapter}{Resume}
Ce rapport presente l'etude technique complete d'un projet informatique de cybersécurité appliquee: une plateforme de simulation de phishing assistee par intelligence artificielle, couplee a un moteur de scoring du risque humain. Le systeme est compose d'un backend expose sous forme d'API REST et d'un client frontend permettant la gestion des campagnes, l'observation comportementale, l'analyse du risque et la remédiation par formation.

La solution couvre l'ensemble du cycle metier: conception et lancement de campagnes, generation de contenus de phishing simules, tracking des interactions (clic, signalement, completion de formation), calcul de scores de risque individuels et consolides, puis pilotage via tableaux de bord. Le present document decrit le contexte, l'architecture, la conception, les choix techniques, l'implementation, la strategie de test, le deploiement et les perspectives d'evolution.

\chapter{Introduction generale du projet}

\section{Contexte du projet}
Les organisations font face a une hausse continue des incidents de cybersécurité lies a l'ingenierie sociale. Les attaques de phishing, en particulier, exploitent la dimension humaine plus que les failles techniques pures. La generalisation du teletravail, des outils cloud et des flux email transverses augmente la probabilite d'exposition.

Dans ce contexte, les programmes de sensibilisation traditionnels montrent leurs limites:
\begin{itemize}[leftmargin=1.2cm]
  \item contenus de formation standardises, peu personnalises;
  \item faible mesure objective des comportements reels;
  \item difficulte a relier les actions de sensibilisation a une reduction mesurable du risque.
\end{itemize}

Le projet vise donc a fournir une plateforme operationnelle qui transforme la sensibilisation en un processus continu, mesurable et améliorable.

\section{Problematique}
Comment concevoir un systeme qui:
\begin{enumerate}[leftmargin=1.2cm]
  \item simule des campagnes de phishing realistes mais controlees;
  \item mesure finement les comportements des utilisateurs;
  \item calcule un niveau de risque humain interpretable;
  \item propose des actions correctives personnalisees;
  \item reste ergonomique pour les equipes metier et securite?
\end{enumerate}

La difficulte principale reside dans l'equilibre entre realisme technique, securite d'usage, explicabilite des scores, et simplicite d'exploitation.

\section{Objectifs}
\subsection{Objectif principal}
Developper une plateforme full-stack backend/frontend permettant de simuler des attaques de phishing de maniere encadree et d'evaluer dynamiquement le risque humain.

\subsection{Objectifs specifiques}
\begin{itemize}[leftmargin=1.2cm]
  \item Mettre en place une API robuste pour la gestion des campagnes et du tracking.
  \item Integrer une generation de contenu email assistee par IA.
  \item Concevoir une interface web de pilotage claire et productive.
  \item Produire des indicateurs analytiques decisionnels.
  \item Intégrer un mecanisme de remédiation post-incident (formation ciblee).
\end{itemize}

\section{Demarche du rapport}
Ce document suit une progression academique:
\begin{itemize}[leftmargin=1.2cm]
  \item de la comprehension du besoin vers la conception;
  \item de la conception vers l'implementation;
  \item de l'implementation vers la validation et le deploiement;
  \item de l'etat actuel vers les perspectives d'industrialisation.
\end{itemize}

\chapter{Presentation globale et architecture du systeme}

\section{Description generale}
Le systeme etudie est une application client-serveur avec:
\begin{itemize}[leftmargin=1.2cm]
  \item un \textbf{backend} centralisant la logique metier, l'authentification, les operations sur les campagnes et l'analytique;
  \item un \textbf{frontend} offrant une interface de gestion des campagnes, d'analyse de risque et de reporting;
  \item une integration \textbf{IA} pour la generation de contenus de phishing et d'articles de formation.
\end{itemize}

Le systeme permet notamment:
\begin{enumerate}[leftmargin=1.2cm]
  \item de definir les cibles (contacts, departements);
  \item de configurer des campagnes (difficulte, type d'attaque);
  \item de generer et envoyer des emails de simulation;
  \item de suivre les interactions des utilisateurs;
  \item de calculer des scores et KPIs;
  \item de declencher des remediations pedagogiques.
\end{enumerate}

\section{Architecture globale client-serveur}
\subsection{Composant client (frontend)}
Le frontend est une SPA (Single Page Application). Il gere:
\begin{itemize}[leftmargin=1.2cm]
  \item l'authentification et la session utilisateur;
  \item la navigation multi-pages metier;
  \item l'affichage des donnees analytiques;
  \item la communication HTTP avec le backend.
\end{itemize}

\subsection{Composant serveur (backend)}
Le backend expose des endpoints REST versionnes, protege les operations sensibles, persiste les donnees en base et encapsule les regles metier (lancement de campagne, mise a jour de score, etc.).

\subsection{Description textuelle du schema d'architecture}
Le diagramme d'architecture globale (non dessine ici) peut etre decrit comme suit:
\begin{itemize}[leftmargin=1.2cm]
  \item un bloc \textit{Client Web} communique en HTTPS vers \textit{API Backend};
  \item le backend lit/ecrit dans la \textit{Base de donnees};
  \item le backend appelle un service externe \textit{IA Generative} pour produire des contenus;
  \item un service mail est utilise pour l'envoi des emails de simulation/formation;
  \item des routes de tracking publiques enregistrent les interactions des destinataires.
\end{itemize}

\section{Choix technologiques}
\subsection{Backend}
\begin{itemize}[leftmargin=1.2cm]
  \item Framework: Laravel (PHP)
  \item Authentification: Laravel Sanctum
  \item ORM: Eloquent
  \item API REST versionnee
\end{itemize}

\subsection{Frontend}
\begin{itemize}[leftmargin=1.2cm]
  \item Framework: React + TypeScript
  \item Routing: React Router
  \item Data fetching/cache: React Query
  \item Client HTTP: Axios
  \item UI: composants modernes type shadcn + Tailwind
\end{itemize}

\subsection{Justification des choix}
Ces technologies ont ete choisies pour:
\begin{itemize}[leftmargin=1.2cm]
  \item la productivite de developpement;
  \item la maturite de l'ecosysteme;
  \item la maintenabilite du code;
  \item la capacite a construire une application metier evolutive.
\end{itemize}

\chapter{Etude et specification des besoins}

\section{Besoins fonctionnels}
\subsection{Gestion des utilisateurs et authentification}
\begin{itemize}[leftmargin=1.2cm]
  \item inscription et connexion;
  \item recuperation du profil utilisateur connecte;
  \item mise a jour du profil;
  \item deconnexion securisee.
\end{itemize}

\subsection{Gestion des contacts cibles}
\begin{itemize}[leftmargin=1.2cm]
  \item creation/modification/suppression de contacts;
  \item import en masse;
  \item classification par departement, poste, seniorite;
  \item consultation des donnees de risque et d'interactions.
\end{itemize}

\subsection{Gestion des campagnes de phishing}
\begin{itemize}[leftmargin=1.2cm]
  \item creation de campagne (nom, difficulte, type d'attaque, ciblage);
  \item generation de contenus;
  \item lancement, pause, reprise, relance;
  \item suivi de statut (draft, scheduled, active, paused, completed).
\end{itemize}

\subsection{Tracking comportemental}
\begin{itemize}[leftmargin=1.2cm]
  \item enregistrement des clics;
  \item enregistrement des signalements;
  \item suivi de completion de formation;
  \item historisation des evenements par contact/campagne.
\end{itemize}

\subsection{Scoring et analytics}
\begin{itemize}[leftmargin=1.2cm]
  \item score de risque par utilisateur;
  \item distribution des risques par niveau;
  \item tendances temporelles;
  \item metriques globales et departementales.
\end{itemize}

\subsection{Formation post-incident}
\begin{itemize}[leftmargin=1.2cm]
  \item attribution automatique d'un module apres clic;
  \item envoi d'un contenu pedagogique contextualise;
  \item suivi de progression et completion.
\end{itemize}

\section{Besoins non fonctionnels}
\subsection{Securite}
\begin{itemize}[leftmargin=1.2cm]
  \item authentification token-based;
  \item validation stricte des entrees API;
  \item separation routes publiques/protegees;
  \item journalisation des erreurs critiques.
\end{itemize}

\subsection{Performance}
\begin{itemize}[leftmargin=1.2cm]
  \item temps de reponse API acceptable en usage nominal;
  \item gestion de volume de contacts/campagnes;
  \item support des operations asynchrones (envoi massif, generation IA).
\end{itemize}

\subsection{Maintenabilite}
\begin{itemize}[leftmargin=1.2cm]
  \item architecture modulaire;
  \item separation claire presentation/metier/persistance;
  \item code lisible et extensible.
\end{itemize}

\subsection{Disponibilite et fiabilite}
\begin{itemize}[leftmargin=1.2cm]
  \item resilience en cas de panne service IA (fallback);
  \item gestion des erreurs d'envoi email;
  \item reprise operationnelle et relance de campagnes.
\end{itemize}

\section{Cas d'utilisation (Use Case Diagram avec explication)}
\subsection{Acteurs}
\begin{itemize}[leftmargin=1.2cm]
  \item \textbf{Administrateur securite}: configure et pilote le systeme.
  \item \textbf{Analyste / Manager}: consulte analytics et rapports.
  \item \textbf{Utilisateur cible}: recoit les emails et interagit (clic/report/formation).
\end{itemize}

\subsection{Description textuelle du diagramme Use Case}
Le diagramme de cas d'utilisation comporte trois acteurs principaux connectes aux cas:
\begin{itemize}[leftmargin=1.2cm]
  \item \textit{Administrateur securite} $\rightarrow$ Gerer contacts, Creer campagne, Lancer campagne, Suspendre campagne, Configurer types d'attaque, Consulter KPI.
  \item \textit{Analyste} $\rightarrow$ Consulter dashboards, Exporter rapports, Analyser tendances.
  \item \textit{Utilisateur cible} $\rightarrow$ Ouvrir email, Signaler phishing, Completer formation.
\end{itemize}

Relations \textit{include/extend} decrites:
\begin{itemize}[leftmargin=1.2cm]
  \item \textit{Lancer campagne} inclut \textit{Generer emails} puis \textit{Envoyer emails}.
  \item \textit{Traiter clic} inclut \textit{Mise a jour score} et \textit{Attribuer formation}.
  \item \textit{Consulter analytics} inclut \textit{Agreger metriques}.
\end{itemize}

\chapter{Conception technique du systeme}

\section{Architecture backend}
\subsection{Couches logiques}
Le backend suit une organisation en couches:
\begin{itemize}[leftmargin=1.2cm]
  \item \textbf{Routes API}: point d'entree HTTP.
  \item \textbf{Controllers}: orchestration des requetes/reponses.
  \item \textbf{Services}: logique metier complexe (PhishingService, RLAgentService).
  \item \textbf{Models ORM}: acces et relations aux donnees.
  \item \textbf{Migrations}: evolution du schema.
\end{itemize}

\subsection{Conception orientee domaine}
Le domaine central est compose de sous-domaines:
\begin{itemize}[leftmargin=1.2cm]
  \item \textit{Identity \& Access};
  \item \textit{Campaign Management};
  \item \textit{Behavior Tracking};
  \item \textit{Risk Scoring};
  \item \textit{Training Remediation};
  \item \textit{Analytics \& Reporting}.
\end{itemize}

\section{Architecture frontend}
\subsection{Organisation fonctionnelle}
Le frontend est organise par pages metier:
\begin{itemize}[leftmargin=1.2cm]
  \item Dashboard
  \item Campaigns
  \item Risk Analytics
  \item Analytics
  \item Training
  \item Users
  \item Reports
  \item Settings / Profile
\end{itemize}

Cette approche favorise la lisibilite fonctionnelle et la separation des responsibilities UI.

\subsection{Navigation et garde d'acces}
La navigation est centralisee dans un routeur applicatif. Une garde \texttt{PrivateRoute} assure que les pages sensibles ne sont accessibles qu'aux sessions authentifiees.

\section{Diagrammes UML (description textuelle)}
\subsection{Diagramme de classes}
Le diagramme de classes inclut principalement:
\begin{itemize}[leftmargin=1.2cm]
  \item \texttt{Campaign} lie a \texttt{CampaignMetrics}, \texttt{SentPhishingEmail}, \texttt{BehavioralEvent};
  \item \texttt{Contact} lie a \texttt{UserRiskScore}, \texttt{Training}, \texttt{BehavioralEvent};
  \item \texttt{RLPolicy} lie a \texttt{Campaign} et potentiellement \texttt{Contact};
  \item \texttt{TrainingModule} lie a \texttt{Training}.
\end{itemize}

Ce modele traduit une relation forte entre simulation (campagne), comportement (events) et remediation (training).

\subsection{Diagramme de sequence: lancement de campagne}
Sequence decrite:
\begin{enumerate}[leftmargin=1.2cm]
  \item L'administrateur soumet ``Lancer campagne'' depuis le frontend.
  \item Le frontend appelle \texttt{POST /campaigns/\{id\}/launch}.
  \item Le controller backend delegue au \texttt{PhishingService}.
  \item Le service recupere les emails \texttt{pending}, envoie via mailer, met a jour les statuts.
  \item Le service met a jour les metriques campagne.
  \item L'API renvoie confirmation et etat actualise.
  \item Le frontend rafraichit la liste des campagnes et affiche un toast.
\end{enumerate}

\subsection{Diagramme de sequence: clic utilisateur}
\begin{enumerate}[leftmargin=1.2cm]
  \item L'utilisateur clique le lien tracke.
  \item Requete publique sur \texttt{/track/click/\{token\}}.
  \item Backend identifie le mail concerne.
  \item Creation d'evenement comportemental ``click''.
  \item Mise a jour score utilisateur.
  \item Eventuelle mise a jour RL reward.
  \item Attribution d'une formation.
  \item Retour reponse de tracking.
\end{enumerate}

\subsection{Diagramme de sequence: signalement}
\begin{enumerate}[leftmargin=1.2cm]
  \item L'utilisateur clique ``Signaler ce phishing''.
  \item API enregistre un evenement ``report''.
  \item Score de risque diminue.
  \item Metriques campagne ajustees.
\end{enumerate}

\section{Modelisation de la base de donnees}
\subsection{Tables centrales}
\begin{longtable}{p{4.3cm} p{8.4cm}}
\toprule
\textbf{Table} & \textbf{Description} \\
\midrule
\endhead
users & Utilisateurs applicatifs (authentification) \\
contacts & Cibles des campagnes (profil RH/metier) \\
campaigns & Campagnes de simulation et configuration \\
sent\_phishing\_emails & Emails generes/envoies/statuts/tracking token \\
behavioral\_events & Evenements clic/signalement avec horodatage \\
user\_risk\_scores & Score et niveau de risque par contact \\
training\_modules & Catalogue de modules pedagogiques \\
trainings & Assignations et progression de formation \\
campaign\_metrics & KPIs de campagne (CTR, precision, AUC) \\
rl\_policies & Donnees d'adaptation comportementale \\
\bottomrule
\end{longtable}

\subsection{Description de l'ERD (sans image)}
L'ERD peut etre lu ainsi:
\begin{itemize}[leftmargin=1.2cm]
  \item \texttt{users (1) --- (N) campaigns}
  \item \texttt{contacts (1) --- (N) behavioral\_events}
  \item \texttt{campaigns (1) --- (N) behavioral\_events}
  \item \texttt{contacts (1) --- (1) user\_risk\_scores}
  \item \texttt{campaigns (1) --- (N) sent\_phishing\_emails}
  \item \texttt{contacts (1) --- (N) trainings}
  \item \texttt{training\_modules (1) --- (N) trainings}
\end{itemize}

\chapter{Implementation du backend (API et logique metier)}

\section{Technologies utilisees}
\begin{itemize}[leftmargin=1.2cm]
  \item PHP 8.x
  \item Laravel
  \item Eloquent ORM
  \item Sanctum
  \item Mail subsystem
  \item Intégration IA via HTTP (Gemini API)
\end{itemize}

\section{Structure du projet backend}
Organisation representative:
\begin{itemize}[leftmargin=1.2cm]
  \item \texttt{app/Http/Controllers/Api}: controllers REST
  \item \texttt{app/Services}: logique metier
  \item \texttt{app/Models}: modeles Eloquent
  \item \texttt{routes/api.php}: contrat API
  \item \texttt{database/migrations}: schema de donnees
\end{itemize}

\section{Endpoints API}
\subsection{Authentification}
\begin{itemize}[leftmargin=1.2cm]
  \item \texttt{POST /api/v1/register}
  \item \texttt{POST /api/v1/login}
  \item \texttt{POST /api/v1/logout}
  \item \texttt{GET /api/v1/user}
  \item \texttt{PUT /api/v1/user}
\end{itemize}

\subsection{Metier principal}
\begin{itemize}[leftmargin=1.2cm]
  \item \texttt{/contacts} (CRUD + import + risk score)
  \item \texttt{/campaigns} (CRUD + pause/resume/launch + metrics)
  \item \texttt{/analytics/*}
  \item \texttt{/trainings, /training-modules}
  \item \texttt{/reports, /email-templates, /departments}
\end{itemize}

\subsection{Tracking public}
\begin{itemize}[leftmargin=1.2cm]
  \item \texttt{/track/click/\{token\}}
  \item \texttt{/track/report/\{token\}}
  \item \texttt{/track/training/\{training\}}
\end{itemize}

\section{Gestion de l'authentification et securite}
\subsection{Authentification Sanctum}
Le backend emet un token personnel a la connexion. Ce token est fourni dans l'en-tete \texttt{Authorization: Bearer ...}. Les routes sensibles sont protegees par le middleware \texttt{auth:sanctum}.

\subsection{Validation des donnees}
Chaque endpoint important valide le payload entrant (types, champs obligatoires, formats, enums). Cette discipline reduit les erreurs applicatives et les vecteurs d'injection.

\subsection{Points de securite appliques}
\begin{itemize}[leftmargin=1.2cm]
  \item hash des mots de passe;
  \item controle d'acces sur endpoints internes;
  \item trace d'erreurs d'envoi et d'appels externes.
\end{itemize}

\subsection{Ameliorations securite recommandees}
\begin{itemize}[leftmargin=1.2cm]
  \item renforcement rate limiting sur routes publiques de tracking;
  \item signature HMAC des tokens/URLs;
  \item durcissement CORS et politique CSP;
  \item journal d'audit enrichi des actions administrateur.
\end{itemize}

\section{Exemples de code expliques}
\subsection{Injection du token de tracking}
Principe: le backend remplace les liens neutres \texttt{href="\#"} par une URL unique contenant un token aléatoire. Cette transformation garantit que chaque interaction est traçable au niveau individu/campagne.

\begin{lstlisting}
// Concept simplifie
$token = Str::random(40);
$trackingUrl = url("/api/v1/track/click/{$token}");
$html = str_replace('href="#"', "href='{$trackingUrl}'", $html);
\end{lstlisting}

Explication:
\begin{itemize}[leftmargin=1.2cm]
  \item un token unique est attribue a chaque e-mail;
  \item le click event devient mesurable;
  \item la campagne peut calculer CTR et patterns comportementaux.
\end{itemize}

\subsection{Mise a jour du score de risque}
\begin{lstlisting}
if ($action === 'click') {
   $newScore = min(100, $newScore + 15);
} elseif ($action === 'report') {
   $newScore = max(0, $newScore - 10);
}
\end{lstlisting}

Explication:
\begin{itemize}[leftmargin=1.2cm]
  \item un clic augmente la vulnerabilite observee;
  \item un signalement traduit une posture defensive;
  \item des bornes [0,100] assurent la stabilite de l'indicateur.
\end{itemize}

\subsection{Adaptation simplifiee type RL}
\begin{lstlisting}
if (rand(1,100) <= 15) {
  // exploration
  return randomAction();
}
// exploitation basee sur click_rate/report_rate
\end{lstlisting}

Explication:
\begin{itemize}[leftmargin=1.2cm]
  \item l'exploration evite la rigidite des scenarios;
  \item l'exploitation personnalise selon comportement historique;
  \item compromis utile entre simplicite et adaptation.
\end{itemize}

\chapter{Implementation du frontend (interface client)}

\section{Technologies utilisees}
\begin{itemize}[leftmargin=1.2cm]
  \item React + TypeScript
  \item React Router
  \item React Query
  \item Axios
  \item Composants UI + Tailwind CSS
\end{itemize}

\section{Structure du projet frontend}
Organisation representative:
\begin{itemize}[leftmargin=1.2cm]
  \item \texttt{src/pages}: pages metier
  \item \texttt{src/components}: composants UI/layout/charts
  \item \texttt{src/services/api.ts}: client HTTP et wrappers d'API
  \item \texttt{src/hooks/useApi.ts}: hooks de requetes/mutations
  \item \texttt{src/store}: etat global session
\end{itemize}

\section{Gestion des etats}
\subsection{Etat de session}
Le store global maintient le token, l'utilisateur courant et le statut d'authentification. Au chargement, l'application tente de recuperer le profil utilisateur; en cas d'echec 401, elle invalide la session.

\subsection{Etat serveur cache}
React Query est utilise pour:
\begin{itemize}[leftmargin=1.2cm]
  \item cacher les resultats d'API;
  \item eviter les requetes redondantes;
  \item invalider intelligemment les donnees apres mutations.
\end{itemize}

\subsection{Etat UI local}
Chaque page gere ses filtres, dialogs, formulaires, modes d'edition et feedback utilisateur (toasts, loading states).

\section{Communication frontend-backend}
\subsection{Client Axios central}
Le client Axios:
\begin{itemize}[leftmargin=1.2cm]
  \item ajoute automatiquement le token;
  \item uniforme les headers;
  \item intercepte les 401 et redirige vers login.
\end{itemize}

\subsection{Pattern appels API}
Les appels sont encapsules par domaine:
\begin{itemize}[leftmargin=1.2cm]
  \item \texttt{authApi}
  \item \texttt{campaignsApi}
  \item \texttt{usersApi}
  \item \texttt{analyticsApi}
  \item \texttt{trainingApi}
\end{itemize}

Ce pattern rend l'integration lisible et facilite l'evolution des endpoints.

\section{Interfaces principales avec explications}
\subsection{Interface ``Campaigns''}
Fonctions majeures:
\begin{itemize}[leftmargin=1.2cm]
  \item creation/edition d'une campagne;
  \item choix type d'attaque et difficulte;
  \item ciblage par departements ou contacts;
  \item lancement/pause/reprise/relance;
  \item visualisation du contenu email;
  \item export PDF/Excel.
\end{itemize}

Valeur ajoutee UX:
\begin{itemize}[leftmargin=1.2cm]
  \item badges de statut;
  \item dialogues de confirmation sur actions irreversibles;
  \item affichage de progression et metriques;
  \item verrouillage en lecture seule pour campagnes deja lancees.
\end{itemize}

\subsection{Interface ``Risk Analytics''}
Fonctions:
\begin{itemize}[leftmargin=1.2cm]
  \item score global de risque;
  \item tendances par departement;
  \item radar de vecteurs d'attaque psychologique;
  \item classement utilisateurs a haut risque.
\end{itemize}

Utilite metier:
\begin{itemize}[leftmargin=1.2cm]
  \item priorisation des actions de formation;
  \item visualisation manageriale des expositions;
  \item suivi des effets des campagnes dans le temps.
\end{itemize}

\subsection{Interface ``Reports''}
Elle consolide les indicateurs pour communication transversale (RSSI, RH, direction) et facilite l'export.

\keypoint{Conclusion de la section frontend}{
L'implementation frontend privilegie la lisibilite operationnelle, la rapidite d'execution des actions metier et la coherence ergonomique. La structure par pages fonctionnelles, combinee a une couche API centralisee, garantit une bonne maintenabilite et facilite les evolutions futures.
}

\chapter{Captures d'ecran et description de l'application}

Cette section est dediee a la documentation visuelle de l'application. Chaque sous-section contient un emplacement de capture, une description fonctionnelle de l'ecran, les actions disponibles et le parcours utilisateur attendu.

\keypoint{Palette visuelle recommandee pour export PDF/Word}{
\textbf{Titres principaux}: bleu profond (\#1F4E79) \\
\textbf{Sous-titres}: bleu secondaire (\#2D6A9F) \\
\textbf{Encadres et zones d'information}: fond gris clair (\#F4F7FA) avec bordure (\#D7DEE8) \\
\textbf{Texte courant}: gris anthracite (\#2E3440)
}

\section{Ecran de connexion}
\captureplaceholder
\textbf{Description de l'interface:} cet ecran constitue le point d'entree de l'application. Il presente un formulaire d'authentification simple avec champs d'identification et actions de connexion.

\textbf{Fonctionnalites disponibles:}
\begin{itemize}[leftmargin=1.2cm]
  \item saisie de l'email et du mot de passe;
  \item validation des champs en cas d'erreur;
  \item redirection automatique vers le dashboard apres succes.
\end{itemize}

\textbf{Interaction utilisateur:} l'utilisateur saisit ses identifiants, soumet le formulaire, puis accede aux pages protegees selon son role.

\section{Dashboard principal}
\captureplaceholder
\textbf{Description de l'interface:} le dashboard centralise les indicateurs critiques (risque global, tendances recentes, campagnes actives). Il offre une vision synthétique pour la prise de decision rapide.

\textbf{Fonctionnalites disponibles:}
\begin{itemize}[leftmargin=1.2cm]
  \item consultation des KPI globaux;
  \item visualisation des tendances de risque;
  \item acces rapide vers les modules Campagnes, Analytics et Rapports.
\end{itemize}

\textbf{Interaction utilisateur:} l'analyste identifie les alertes prioritaires puis navigue vers la page la plus pertinente pour investiguer.

\section{Page de gestion des campagnes}
\captureplaceholder
\textbf{Description de l'interface:} cette page est le centre operationnel de la plateforme. Elle affiche la liste des campagnes avec statuts, difficultes, metriques et actions contextuelles.

\textbf{Fonctionnalites disponibles:}
\begin{itemize}[leftmargin=1.2cm]
  \item creation, modification et suppression de campagnes;
  \item filtres par statut, difficulte et recherche textuelle;
  \item lancement, pause, reprise et relance des campagnes.
\end{itemize}

\textbf{Interaction utilisateur:} l'administrateur planifie une campagne, verifie sa configuration, puis la lance et suit son evolution en temps reel.

\section{Formulaire de creation de campagne}
\captureplaceholder
\textbf{Description de l'interface:} ce formulaire regroupe les paramètres strategiques de simulation: nom, type d'attaque, niveau de difficulte, ciblage et options d'automatisation.

\textbf{Fonctionnalites disponibles:}
\begin{itemize}[leftmargin=1.2cm]
  \item selection du type d'attaque et du niveau de difficulte;
  \item ciblage global, par departement ou par utilisateur;
  \item previsualisation du contenu e-mail avant validation.
\end{itemize}

\textbf{Interaction utilisateur:} l'operateur ajuste les paramètres puis enregistre la campagne. Un retour visuel confirme la prise en compte.

\section{Analyse de risque (Risk Analytics)}
\captureplaceholder
\textbf{Description de l'interface:} cet ecran fournit une vue analytique approfondie du risque humain, avec graphiques de tendance, radar de vulnerabilites et classement des profils sensibles.

\textbf{Fonctionnalites disponibles:}
\begin{itemize}[leftmargin=1.2cm]
  \item filtre par departement;
  \item tableau des utilisateurs a haut risque;
  \item export des analyses vers PDF et Excel.
\end{itemize}

\textbf{Interaction utilisateur:} le responsable securite identifie les zones de risque eleve, priorise les actions de remédiation et suit l'impact des campagnes.

\section{Page Utilisateurs / Contacts}
\captureplaceholder
\textbf{Description de l'interface:} la page utilisateurs permet de consulter le referentiel des cibles, leurs informations metier et leur historique comportemental.

\textbf{Fonctionnalites disponibles:}
\begin{itemize}[leftmargin=1.2cm]
  \item ajout et edition de profils;
  \item import en masse de contacts;
  \item consultation du score individuel et de l'historique.
\end{itemize}

\textbf{Interaction utilisateur:} l'administrateur met a jour le portefeuille de cibles et segmente les populations pour les futures campagnes.

\section{Page Rapports}
\captureplaceholder
\textbf{Description de l'interface:} cet ecran regroupe les resultats consolidés des campagnes et facilite la communication vers les parties prenantes.

\textbf{Fonctionnalites disponibles:}
\begin{itemize}[leftmargin=1.2cm]
  \item consultation des indicateurs consolides;
  \item generation et telechargement de rapports;
  \item comparaison inter-campagnes.
\end{itemize}

\textbf{Interaction utilisateur:} l'analyste prepare un rapport de synthese pour le management et documente les actions recommandees.

\section{Page Formation}
\captureplaceholder
\textbf{Description de l'interface:} la section formation presente les modules assignes, l'etat de progression et les contenus pedagogiques post-incident.

\textbf{Fonctionnalites disponibles:}
\begin{itemize}[leftmargin=1.2cm]
  \item suivi des assignations et completions;
  \item consultation des contenus de remédiation;
  \item priorisation des utilisateurs en retard de formation.
\end{itemize}

\textbf{Interaction utilisateur:} le responsable programme suit les taux de completion et ajuste les campagnes de sensibilisation.

\section{Parametres et profil}
\captureplaceholder
\textbf{Description de l'interface:} les ecrans ``Profil'' et ``Parametres'' permettent d'adapter les informations de compte et les preferances d'usage.

\textbf{Fonctionnalites disponibles:}
\begin{itemize}[leftmargin=1.2cm]
  \item modification des informations personnelles;
  \item mise a jour des preferances organisationnelles;
  \item gestion des parametres globaux de l'application.
\end{itemize}

\textbf{Interaction utilisateur:} chaque utilisateur ajuste son espace de travail pour une experience plus fluide et conforme a son contexte.

\keypoint{Conclusion de la section Captures d'ecran}{
L'ajout de captures commentees renforce la valeur pedagogique du rapport. Cette section facilite la comprehension de l'application par un jury, un encadrant ou une equipe de reprise, en reliant directement l'architecture technique aux ecrans reels d'exploitation.
}

\chapter{Strategie de tests et validation}

\section{Types de tests}
\subsection{Tests unitaires}
Objectif: verifier les unites de logique metier isolees (calcul score, transitions de statut, selection de cibles).

Exemples:
\begin{itemize}[leftmargin=1.2cm]
  \item test de variation score click/report;
  \item test de classification des niveaux de risque;
  \item test de remplacement correct des liens de tracking.
\end{itemize}

\subsection{Tests d'integration}
Objectif: verifier l'interaction entre composants backend (controller + service + ORM + DB).

Exemples:
\begin{itemize}[leftmargin=1.2cm]
  \item creation campagne puis generation emails;
  \item route tracking click puis creation event et mise a jour score;
  \item completion de training puis impact analytics.
\end{itemize}

\subsection{Tests end-to-end (E2E)}
Objectif: valider les parcours utilisateurs complets frontend+backend.

Scenarios:
\begin{itemize}[leftmargin=1.2cm]
  \item login $\rightarrow$ creer campagne $\rightarrow$ launch $\rightarrow$ verifier statut;
  \item consulter analytics puis exporter rapport;
  \item ouvrir profil utilisateur et verifier historique de risque.
\end{itemize}

\subsection{Tests de non-regression}
Apres chaque evolution majeure:
\begin{itemize}[leftmargin=1.2cm]
  \item verification du contrat API;
  \item verification du rendu des pages critiques;
  \item verification des cas limites (campagne sans cibles, token invalide, etc.).
\end{itemize}

\section{Outils utilises}
\begin{itemize}[leftmargin=1.2cm]
  \item Backend: PHPUnit/Laravel test runner (recommande)
  \item Frontend: Vitest/Jest + Testing Library (recommande)
  \item E2E: Playwright/Cypress (recommande)
  \item Qualite statique: ESLint + type checking TypeScript
\end{itemize}

\section{Plan de validation propose}
\begin{longtable}{p{5cm} p{4cm} p{5cm}}
\toprule
\textbf{Objet teste} & \textbf{Type de test} & \textbf{Critere de succes} \\
\midrule
\endhead
Authentification & Integration & Login renvoie token + user, routes protegees inaccessibles sans token \\
Creation campagne & Integration & Donnees persistantes, metriques initialisees, pas d'erreur validation \\
Launch campagne & Integration & Emails statuts ``sent'' ou ``failed'', campagne active \\
Tracking click & Integration & Event cree, score augmente, formation assignee \\
Tracking report & Integration & Event cree, score diminue, metriques ajustees \\
Liste campagnes frontend & E2E & Affichage coherent avec filtres et statuts \\
Analyse de risque & E2E & Graphes et tableaux charges sans erreur API \\
Exports & E2E & Fichier genere et telecharge correctement \\
\bottomrule
\end{longtable}

\section{Resultats des tests}
Dans l'etat actuel du projet, l'infrastructure de tests automatisee est encore perfectible. Toutefois, les validations manuelles des parcours majeurs montrent:
\begin{itemize}[leftmargin=1.2cm]
  \item bonne stabilite fonctionnelle des flux principaux;
  \item cohérence globale des interactions campagne/tracking/scoring;
  \item besoin de renforcer l'automatisation pour fiabiliser les evolutions futures.
\end{itemize}

\chapter{Deploiement et mise en production}

\section{Environnement de deploiement}
\subsection{Developpement}
\begin{itemize}[leftmargin=1.2cm]
  \item Backend lance localement (serveur PHP/Laravel);
  \item Frontend lance via Vite en mode dev;
  \item Base de donnees locale de test.
\end{itemize}

\subsection{Pre-production}
Environnement miroir plus proche de la production:
\begin{itemize}[leftmargin=1.2cm]
  \item variables d'environnement separees;
  \item service mail de test;
  \item monitoring de base.
\end{itemize}

\subsection{Production}
Architecture cible recommandee:
\begin{itemize}[leftmargin=1.2cm]
  \item serveur applicatif backend;
  \item stockage DB gere;
  \item frontend statique servi par CDN/reverse proxy;
  \item service de queue pour traitements lourds;
  \item supervision logs + metriques.
\end{itemize}

\section{Outils et plateformes}
\begin{itemize}[leftmargin=1.2cm]
  \item CI/CD (GitHub Actions/GitLab CI)
  \item Conteneurisation Docker (option recommandee)
  \item Reverse proxy Nginx
  \item Monitoring (Grafana/Prometheus ou equivalent)
  \item Gestion des secrets (vault ou variables chiffrees)
\end{itemize}

\section{Procedure de mise en production}
\subsection{Etapes techniques}
\begin{enumerate}[leftmargin=1.2cm]
  \item validation branche release;
  \item execution tests automatises;
  \item build frontend;
  \item migration base de donnees;
  \item deploiement backend;
  \item deploiement assets frontend;
  \item verification smoke tests;
  \item ouverture progressive trafic.
\end{enumerate}

\subsection{Checklist de go-live}
\begin{itemize}[leftmargin=1.2cm]
  \item variables \texttt{APP\_KEY}, DB, MAIL, IA configurees;
  \item HTTPS actif;
  \item politique de sauvegarde active;
  \item plan rollback documente;
  \item alerting operationnel configure.
\end{itemize}

\section{Plan de rollback}
En cas d'incident post-deploiement:
\begin{enumerate}[leftmargin=1.2cm]
  \item retour a la version precedente de l'application;
  \item restauration schema si migration critique;
  \item verification integrite donnees;
  \item communication interne incident.
\end{enumerate}

\chapter{Difficultes techniques rencontrees et solutions}

\section{Problemes techniques}
\subsection{Complexite du workflow campagne}
Le cycle creation/generation/envoi/tracking/remediation implique plusieurs transitions de statut et de nombreux cas limites (emails deja envoyes, relance, interactions partielles).

\textbf{Solution}: centraliser cette logique dans un service metier unique et introduire des statuts explicites.

\subsection{Integration IA et robustesse}
Les appels IA peuvent echouer (quota, latence, indisponibilite).

\textbf{Solution}: fallback local de generation, retries, journalisation detaillee.

\subsection{Coherence frontend-backend}
Dans une architecture evolutive, certains endpoints peuvent diverger entre implementation backend et wrappers frontend.

\textbf{Solution}: formaliser un contrat API unique (OpenAPI) et le valider en CI.

\subsection{Qualite analytique}
Certaines visualisations de demo utilisent des composantes pseudo-aleatoires pour enrichir l'affichage en phase prototype.

\textbf{Solution}: basculer progressivement vers des indicateurs 100\% derives des donnees tracees.

\section{Solutions apportees}
\begin{itemize}[leftmargin=1.2cm]
  \item decoupage clair des couches et services;
  \item validation stricte des payloads;
  \item ajout de metadonnees operationnelles (statuts, horodatages);
  \item UX orientee prevention d'erreurs (dialogs, confirmations, modes read-only).
\end{itemize}

\section{Retours d'experience}
Le principal enseignement est que la valeur d'une plateforme de sensibilisation cyber ne vient pas seulement de la generation de scenarios, mais de la qualite de la boucle de feedback: observation comportementale fiable, scoring interpretable, remediation rapide et pilotage continu.

\chapter{Conclusion generale et perspectives}

\section{Bilan du projet}
Le projet aboutit a une plateforme credible et fonctionnelle couvrant l'essentiel d'un programme de phishing simulation moderne:
\begin{itemize}[leftmargin=1.2cm]
  \item gestion complete du cycle campagne;
  \item tracking comportemental exploitable;
  \item scoring du risque utilisateur/departement;
  \item capacites analytiques orientées decision;
  \item remédiation pedagogique post-incident.
\end{itemize}

L'architecture technique montre une base solide et extensible, avec une bonne separation des responsabilites entre backend et frontend.

\section{Ameliorations futures}
\subsection{Axe technique}
\begin{itemize}[leftmargin=1.2cm]
  \item renforcer couverture de tests auto (unitaires, integration, E2E);
  \item industrialiser le monitoring et l'observabilite;
  \item normaliser strictement les contrats API.
\end{itemize}

\subsection{Axe data/IA}
\begin{itemize}[leftmargin=1.2cm]
  \item enrichir le modele de risque (features temporelles, profils metiers);
  \item introduire des bandits contextuels plus formels;
  \item versionner prompts et evaluer qualite IA.
\end{itemize}

\subsection{Axe metier/gouvernance}
\begin{itemize}[leftmargin=1.2cm]
  \item workflows d'approbation des campagnes sensibles;
  \item rapports executifs automatiques;
  \item alignement complet RGPD et auditabilite.
\end{itemize}

\section{Ouverture}
La plateforme peut evoluer vers un veritable \textit{cockpit de resilience humaine} en cybersécurité, couvrant non seulement l'email phishing mais aussi les vecteurs SMS, messagerie instantanee et voice phishing, avec une intelligence adaptative explicable et un pilotage multi-entites.

\appendix

\chapter{Annexe A - Inventaire des endpoints principaux}
\begin{longtable}{p{5.3cm} p{1.8cm} p{7cm}}
\toprule
\textbf{Endpoint} & \textbf{Methode} & \textbf{Role} \\
\midrule
\endhead
/api/v1/register & POST & Inscription utilisateur \\
/api/v1/login & POST & Connexion \\
/api/v1/logout & POST & Deconnexion \\
/api/v1/user & GET/PUT & Lecture/mise a jour profil \\
/api/v1/contacts & CRUD & Gestion des contacts \\
/api/v1/contacts/import & POST & Import contacts \\
/api/v1/campaigns & CRUD & Gestion campagnes \\
/api/v1/campaigns/\{id\}/launch & POST & Lancer campagne \\
/api/v1/campaigns/\{id\}/pause & POST & Pause campagne \\
/api/v1/campaigns/\{id\}/resume & POST & Reprise campagne \\
/api/v1/campaigns/\{id\}/metrics & GET & Metriques campagne \\
/api/v1/behavioral-events & GET/POST & Historique/create event \\
/api/v1/risk-scores & GET/PUT & Scores de risque \\
/api/v1/trainings & CRUD partiel & Suivi formations \\
/api/v1/training-modules & CRUD & Catalogue formations \\
/api/v1/reports & CRUD & Rapports \\
/api/v1/analytics/click-rate-trend & GET & Tendance click/report \\
/api/v1/analytics/risk-distribution & GET & Distribution risques \\
/api/v1/analytics/campaign-performance & GET & Comparatif campagnes \\
/api/v1/analytics/behavior-heatmap & GET & Heatmap comportementale \\
/api/v1/analytics/training-effectiveness & GET & Efficacite formation \\
/api/v1/analytics/department-risk & GET & Risque par departement \\
/api/v1/analytics/global-metrics & GET & KPIs globaux \\
/api/v1/ai/generate-template & POST & Generation email IA \\
/api/v1/track/click/\{token\} & GET & Tracking clic \\
/api/v1/track/report/\{token\} & GET & Tracking signalement \\
/api/v1/track/training/\{training\} & GET & Validation formation \\
\bottomrule
\end{longtable}

\chapter{Annexe B - Scenarios de tests detailles}
\section*{Scenario 1: Creation et lancement campagne}
\textbf{Precondition}: utilisateur authentifie, contacts presents.

\textbf{Etapes}:
\begin{enumerate}[leftmargin=1.2cm]
  \item Ouvrir page Campagnes.
  \item Creer campagne ``Spear Phishing Q2''.
  \item Choisir difficulte ``difficile''.
  \item Cibler departement ``Finance''.
  \item Enregistrer puis lancer.
\end{enumerate}

\textbf{Resultat attendu}:
\begin{itemize}[leftmargin=1.2cm]
  \item statut campagne $\rightarrow$ active;
  \item e-mails generes et marques sent/pending selon traitement;
  \item metriques initiales visibles.
\end{itemize}

\section*{Scenario 2: Clic utilisateur}
\textbf{Etapes}:
\begin{enumerate}[leftmargin=1.2cm]
  \item Ouvrir email simule.
  \item Cliquer CTA.
  \item Verifier en base: status clicked + behavioral event.
  \item Verifier score utilisateur augmente.
\end{enumerate}

\textbf{Resultat attendu}:
\begin{itemize}[leftmargin=1.2cm]
  \item event click enregistre;
  \item score mis a jour;
  \item formation assignee.
\end{itemize}

\section*{Scenario 3: Signalement}
\textbf{Etapes}:
\begin{enumerate}[leftmargin=1.2cm]
  \item Cliquer ``Signaler ce phishing''.
  \item Verifier event report.
  \item Verifier diminution du score.
\end{enumerate}

\textbf{Resultat attendu}:
\begin{itemize}[leftmargin=1.2cm]
  \item event report enregistre;
  \item score ajuste a la baisse;
  \item metriques campagne rafraichies.
\end{itemize}

\chapter{Annexe C - Guide de lecture des diagrammes UML}
\section*{Use Case}
Observer:
\begin{itemize}[leftmargin=1.2cm]
  \item acteurs externes;
  \item cas metier centraux;
  \item dependances include/extend.
\end{itemize}

\section*{Classes}
Observer:
\begin{itemize}[leftmargin=1.2cm]
  \item entites fortes (Campaign, Contact, UserRiskScore);
  \item cardinalites 1-1 et 1-N;
  \item agregations fonctionnelles.
\end{itemize}

\section*{Sequence}
Observer:
\begin{itemize}[leftmargin=1.2cm]
  \item ordre temporel des interactions;
  \item appels synchrones/asynchrones;
  \item points de persistance et mises a jour d'etat.
\end{itemize}

\chapter{Annexe D - Plan d'amelioration sur 12 mois}
\begin{longtable}{p{2.5cm} p{4.5cm} p{5.5cm} p{2.5cm}}
\toprule
\textbf{Periode} & \textbf{Axe} & \textbf{Actions} & \textbf{Indicateur} \\
\midrule
\endhead
M1-M3 & Qualite & Contrat API OpenAPI, tests unitaires critiques & Couverture > 50\% \\
M3-M6 & Data & Analytics 100\% data-driven & 0 KPI simule \\
M6-M9 & Securite & Hardening tokens tracking, audit log & 100\% actions tracees \\
M9-M12 & IA & Personalisation avancée (bandits contextuels) & +X\% detection precoce \\
\bottomrule
\end{longtable}

\chapter*{Bibliographie indicative}
\addcontentsline{toc}{chapter}{Bibliographie indicative}
\begin{enumerate}[leftmargin=1.2cm]
  \item OWASP Foundation, \textit{OWASP Top 10}.
  \item NIST, \textit{Cybersecurity Awareness and Training Guidance}.
  \item ENISA, \textit{Cybersecurity Culture Guidelines}.
  \item Documentation Laravel, \url{https://laravel.com/docs}.
  \item Documentation React, \url{https://react.dev}.
  \item Documentation TypeScript, \url{https://www.typescriptlang.org/docs}.
  \item Documentation Axios et React Query.
\end{enumerate}

\end{document}
